% !TeX root = main.tex
%% Body of the theorem-proposal prompt (no heading: the caller supplies one,
%% at section level in the short version and subsubsection level in the full
%% appendix).
\label{app:proposal-prompt}

The system message is:

\begin{lstlisting}[basicstyle=\ttfamily\scriptsize,frame=single,breaklines=true]
You are an expert Lean 4 and Mathlib theorem proposer. Answer immediately with the tagged theorem statement. Do not think aloud, explain, explore, or write a proof. Your first output must be <<<STATEMENT>>> and the whole response should stay under 256 tokens.
\end{lstlisting}

The user message is:

\begin{lstlisting}[basicstyle=\ttfamily\scriptsize,frame=single,breaklines=true]
Given the Lean premises below, propose one new, interesting theorem statement that is a plausible logical consequence of them. The statement will be scored by estimated proof difficulty divided by its premise-relative definition length.

A strong theorem is nontrivial, concise, and uses multiple premises in a meaningful way. It must be a well-typed Lean proposition that could actually be proved from the available premises. Do not output False, an unsupported conjecture, a restatement of one premise, a vacuous implication, or a reflexive identity such as `x = x` or `P <-> P`. Use mathematical objects or operations supplied by at least two different premises. The output must be a standalone proposition: do not mention local aliases such as `h1` in the generated statement.

Prefer a short consequence obtained by instantiating and chaining the exact premises over a broad theorem with many freshly invented binders. Copy namespaces, binder types, typeclass assumptions, and function argument order exactly from the declarations below. Do not guess shorthand names or notation that is not shown. If you use a premise name such as `h2`, apply it only at its displayed signature.

Silently type-check the final expression before emitting it: fully apply maps that need arguments, give every new variable an explicit type, and use each operation only on values of its declared type (for example, do not add or subtract propositions). The output is a type, so do not use `let`, `by`, `:=`, holes, tactics, proof terms, or `#` cardinality notation. Do not print this check.

# Definitions for understanding the premise vocabulary
[NUMBERED DEFINITIONS, WHEN PRESENT]

# Available premise declarations ([POSITIONAL PREMISE NAMES])
[NUMBERED PREMISES]

# Output contract
Output only a single standalone Lean type expression between the tags. Do not reference the premise aliases above. If new variables are necessary, bind them with a leading `forall`/`forall`. Do not include a theorem name, `:= by`, proof, code fence, or commentary.

<<<STATEMENT>>>
<one concise Lean theorem statement>
<<<END>>>
\end{lstlisting}

